\section{Represented transport and absolute period errors}
\label{sec:transport}

We now construct the geometric observations used in
Section~\ref{sec:integral-recovery}. The domain is the known ball
$B_{\rm ev}$ of Theorem~\ref{thm:evaluated-periods}. The reference
point $p_*$ has algebraic fibre equations, and
Theorem~\ref{thm:smooth-marking} supplies its marking by smooth
integral chains.

\begin{theorem}\label{thm:accessible-observations}
For every named $p\in B_{\rm ev}^{\circ}$ one can construct a
geometric symplectic marking of $X_p$ by evaluated smooth chains,
and eight evaluated smooth based loops, each a positive meridian
of one discriminant hyperplane and with zero winding about the others.
The initial and endpoint periods of every marked chain along these
loops have arbitrarily accurate absolute enclosures obtained by
finite computation relative to the name of $p$.

The construction uses the fibre equations, their local residue form,
and the independently constructed marking. It uses no local-sphere
coordinate column, intersection vector, endpoint homology
decomposition, or monodromy matrix.
\end{theorem}

\subsection{Paths with known margins}

Let $x,y\in B_{\rm ev}^{\circ}$ be named points. Enumerate
$q\in R\cap(\QQ+\sqrt{-1}\QQ)^8\cap B_{\rm ev}$ until the
segments $[x,q]$ and $[q,y]$ avoid every coordinate hyperplane.
Such $q$ form a dense open set. For a fixed nonzero $x_i$, the
segment from $x_i$ to $q_i$ contains zero precisely when $q_i$
belongs to the closed opposite ray $\{-a x_i:a\ge0\}$. Its inverse
image under the nonzero complex-linear coordinate functional on
$R$ has empty interior. The finitely many conditions for $x$ and
$y$ therefore exclude only a closed nowhere-dense set.

The search is effective by dovetailing candidates and refinements.
At a candidate, a slope bound for each scalar segment comes from
endpoint enclosures. Values on a finite time mesh, less evaluation
errors and the slope bound times the cell diameter, give lower
bounds for its modulus. If the segment avoids zero, sufficiently
fine meshes prove a positive rational lower bound. Bad candidates
need not be decided. Also find positive lower bounds for
$\rho-\|x\|,\rho-\|y\|,\rho-\|q\|$. Convexity of the ball gives
their minimum as an interior margin for the polygonal path.

Fix once the smooth increasing function
\begin{equation}\label{eq:flat-time-change}
 \theta(t)=\frac{\int_0^t b(u)\,\dd u}{\int_0^1 b(u)\,\dd u},
 \qquad
 b(u)=
 \begin{cases}
  \exp[-1/(u(1-u))],&0<u<1,\\
  0,&\text{otherwise}.
 \end{cases}
\end{equation}
Extend $\theta$ by zero and one outside the interval. Every
positive-order derivative vanishes at its endpoints. Values and
bounds for any fixed number of derivatives are computable:
inside the interval the derivatives are rational functions times
the exponential; exponential decay controls their endpoint tails,
and compact interior intervals have elementary bounds. Quadrature
evaluates the integral, whose denominator has a positive lower
bound from any fixed interior subinterval.

Parametrize each segment with $\theta$ and concatenate. This
gives a smooth path with known derivative bounds and the same
positive margins. In particular it constructs a path from $p_*$
to $p$. Its choice need not be canonical.

\subsection{Individual meridians and flat local spheres}

For node $i$, choose $j\ne i$ in the same block and put
$v_i=e_i-e_j\in R$. Choose a rational complex point $a_i$ with
$(a_i)_i=0$, all other coordinates nonzero, and $\|a_i\|<\rho/4$.
Such points exist because no other coordinate hyperplane contains
$\{p_i=0\}\subset R$. A sufficiently small rational $r_i>0$
satisfies
\begin{equation}\label{eq:meridian-disc-bounds}
 \|a_i\|+r_i\|v_i\|<\rho/2,\qquad
 r_i|(v_i)_k|<\tfrac12|(a_i)_k|\quad(k\ne i).
\end{equation}
The disc $a_i+\zeta v_i$, $|\zeta|\le r_i$, meets just the
$i$th hyperplane, transversely at its centre. Its boundary
$a_i+r_i e^{2\pi\sqrt{-1}t}v_i$ has winding one about that
hyperplane and zero about the others. Join its starting point to
$p$ by the preceding path construction, follow the circle, and
return along the reverse path. Flat reparametrization at the joins
gives the required smooth based loop $\ell_i$, with known
derivative bounds, a positive lower bound for every $|p_k|$,
and a known margin inside the ball.

The local Milnor homology at a threefold node is generated by the
oriented real $S^3$ in $\sum w_j^2=\tau_i$. One circuit in
$\tau_i$ changes $w$ to $-w$, whose action on $S^3$ has degree
$+1$. Thus the generator defines a single-valued flat local
section. Inclusion into the compact fibre commutes with transport,
so the classes $\delta_i$ are globally flat on $B_{\rm ev}^{\circ}$.
Their representatives occupy disjoint Milnor balls and have zero
pairwise intersections.

On the transverse disc above,
$\tau_i=\zeta/U_i(a_i+\zeta v_i)$ has a simple zero and winding
$+1$. The total space over the disc is smooth and its central
fibre has just one ordinary double point. This is precisely the
one-critical-point hypothesis of the positive dimension-three
Picard--Lefschetz formula
\cite[Expos\'e~XIV, \S3.2.3 and Theorem~3.2.11]{Deligne}.
Changing the whisker conjugates its action by a symplectic
monodromy fixing $\delta_i$, and hence leaves that action unchanged.
The individual observation uses the full six-dimensional base:
varying $p_i$ alone would violate its block relation.

\subsection{A horizontal field determined by the equations}

Compose a constructed path with $\sigma\pi_R$. Its known compact
domain margin and the bounded analytic presentation of $\sigma$
give evaluations and bounds for every fixed number of derivatives
of the defining sections. Use the fixed Fubini--Study metrics and
connections on $A$ and on the equation bundle
$\mathcal E=\OO(1,1)^{\oplus2}\oplus\OO(2,2)$.
Regard this as a real rank-six bundle for adjoints. If $F(t)$
denotes the three sections and $A_F=\nabla_xF$, then
\begin{equation}\label{eq:ambient-positivity}
 g(t,x)=|F(t,x)|^2+\lambda_{\min}(A_FA_F^*)>0
                    \qquad ((t,x)\in[0,1]\times A).
\end{equation}
A zero would be a common zero with deficient equation derivative,
contrary to smoothness. Compactness makes the lower bound positive.

A rational $m>0$ with $g\ge m$ is found without solving for the
fibre points. On finitely many compact ambient coordinate cells,
evaluate $F,A_FA_F^*$ and their variation bounds. The inequality
$|\lambda_{\min}(B)-\lambda_{\min}(C)|\le\|B-C\|$ controls
the least eigenvalue between mesh points. Refinement of the mesh
and evaluation tolerances eventually proves a positive common lower
bound. This is a strict positive-margin search on the entire compact
ambient space.

On $|F|^2\le m/2$ we have $A_FA_F^*\ge m/2$. Choose a smooth
cutoff $\chi$ equal to one on $[0,m/8]$ and zero on $[m/4,\infty)$.
The field
\begin{equation}\label{eq:equation-horizontal-field}
 V=-\chi(|F|^2)A_F^*(A_FA_F^*)^{-1}\nabla_tF
\end{equation}
on $|F|^2<m/2$, extended by zero elsewhere, is smooth because
its support has a buffer inside the invertibility region. On the
zero set $\nabla_tF+A_FV=0$. Thus $(\partial_t,V)$ is tangent
to the total zero manifold, and its flow transports one actual
fibre to the next. Compactness gives existence for the whole
interval, and reverse transport gives the inverse.

Every derivative bound follows from the equation bounds and $m$.
On the support the inverse has norm at most $2/m$, and
\[
                         D(B^{-1})=-B^{-1}(DB)B^{-1}.
\]
Three derivatives of $F$
suffice for the second spatial derivative of $V$. Higher jet
evaluation below uses correspondingly higher derivatives, supplied
by the bounded analytic presentation.

Embed $A$ by its two rank-one Hermitian projectors in a Euclidean
space. A fixed tubular retraction is obtained near each projector
by taking the separated top eigenspace projector. Its uniform
eigenvalue gap gives derivative bounds. Compose $V$ with this
retraction and a compact cutoff equal to one near $A$, to obtain
a globally defined Euclidean field $U(t,x)$ with bounded derivatives.
It agrees with $V$ on $A$. Its solutions starting there stay there
by uniqueness.

\subsection{Transported derivatives and their evaluation}

For an initial simplex map $f$ write
$x(t,s)=\Phi_t(f(s))$, $P=D_sx$, and $Q=D_s^2x$.
They satisfy
\begin{equation}\label{eq:transport-jets}
 \dot x=U(t,x),\qquad
 \dot P=D_xU(t,x)P,\qquad
 \dot Q=D_xU(t,x)Q+D_x^2U(t,x)[P,P].
\end{equation}
Let $L$ bound $\|D_xU\|$, $M_2$ bound $\|D_x^2U\|$, and
$F_1,F_2$ bound the corresponding initial simplex derivatives on
their given neighbourhoods. Then
\begin{equation}\label{eq:jet-size-bounds}
 \|P\|\le B_1=e^L\max(1,F_1),\qquad
 \|Q\|\le B_2=e^L(F_2+M_2B_1^2).
\end{equation}
The first estimate is Gronwall's inequality; the second follows
by variation of constants in the last equation.

These equations evaluate both derivatives, not just the map. On
$\|P\|\le B_1+1$, $\|Q\|\le B_2+1$, their right side has a
known Lipschitz bound, using $D_x^3U$ in the last component.
Append smooth cutoffs in $P,Q$ which are one there and vanish
outside larger bounded sets. The resulting augmented field
$\mathcal F$ is globally bounded and Lipschitz, and the true
jet solution stays where the cutoffs equal one.

For completeness, suppose $\|\mathcal F\|\le M$,
$\|D_z\mathcal F\|\le L_*$ and
$\|\partial_t\mathcal F\|\le M_t$, enlarging $L_*$ to be at least
one. Its exact solutions have $\|z''\|\le M_t+L_*M$.
Euler steps of length $h=1/N$, with evaluated field errors at
most $\epsilon$ and initial error at most $e_0$, obey
\begin{equation}\label{eq:euler-jet-enclosure}
 \|z_N-z(1)\|\le e^{L_*}e_0+
       \frac{e^{L_*}-1}{L_*}
           \left(\epsilon+\tfrac12(M_t+L_*M)h\right).
\end{equation}
This follows by summing the one-step defect and using
$(1+hL_*)^N\le e^{L_*}$. Refining $N$, $\epsilon$ and $e_0$
therefore evaluates the transported map and its two derivatives
at every named simplex point to any requested accuracy.

The actual maps remain smooth on the original simplex
neighbourhoods, with bounds \eqref{eq:jet-size-bounds}.
They have exact matching face maps because the same actual flow
is applied to the same initial face maps. Their numerical
enclosures are not asserted to constitute a separate exactly
closed chain. Applying this transport from $p_*$ to $p$ gives
the geometric symplectic marking at the target fibre.

\subsection{A first-jet error estimate for integration}

Fix one target cycle
$c=\sum_\nu m_\nu(f_\nu)_\#[\Delta^3]$ and a meridian at $p$.
Let $\bar\Psi_p$ be a smooth ambient extension of the three-form
on the endpoint fibre. Such an extension with evaluated bounds
can be constructed from regular residue charts. Choose nonzero
three-equation minors on a finite cover of the smooth fibre;
the compact positivity bound \eqref{eq:ambient-positivity},
strict minor inequalities and refined coordinate cells supply
an effective cover and denominator bounds. Local residue
expressions extend near those cells. A computable smooth
partition of unity and the ambient tube extend the form to a
compactly supported Euclidean form agreeing on tangent vectors
to $X_p$. This yields known bounds
$W_0=\sup\|\bar\Psi_p\|$, $W_1=\sup\|D\bar\Psi_p\|$.

Use common $F_1,F_2$ for the finitely many simplex maps, and
the constants in \eqref{eq:jet-size-bounds}. Put
\begin{equation}\label{eq:quadrature-constants}
 \begin{split}
 S&=\sum_\nu|m_\nu|\operatorname{vol}(\Delta^3),\\
 C_q&=W_1B_1^4+3W_0B_1^2B_2,\\
 C_j&=W_1(B_1+1)^3+3W_0(B_1+1)^2,\\
 \Lambda_*&=1+2L+M_2(B_1+1).
 \end{split}
\end{equation}
All norms are Euclidean operator norms, with absolute values for
complex coefficients. A zero chain needs no numerical integration.

\begin{proposition}\label{prop:absolute-period-error}
At each quadrature point suppose a first-jet approximation
$\widehat z=(\widehat x,\widehat P)$ is absolutely continuous
and piecewise $C^1$, has initial error at most $e_0$, satisfies
$\|\widehat P\|\le B_1+1$, and has differential-equation residual
at most $r$ in the block maximum norm. Define
\[
 E=e^{\Lambda_*}e_0+
             \frac{e^{\Lambda_*}-1}{\Lambda_*}r<1.
\]
For a simplex mesh of maximal diameter $h$ and scalar
integrand-evaluation error at most $\xi$, the resulting endpoint
sum $\widehat I$ satisfies
\begin{equation}\label{eq:absolute-period-error}
 \left|\widehat I-\int_{(\Phi_1)_\#c}\Psi_p\right|
                      \le S(C_qh+C_jE+\xi).
\end{equation}
\end{proposition}

\begin{proof}
For $z=(x,P)$ the right side of the first two jet equations is
\[
             \mathcal F_1(t,z)=(U(t,x),D_xU(t,x)P).
\]
On the convex set
$\|P\|\le B_1+1$ it has Lipschitz constant at most
$L+M_2(B_1+1)\le\Lambda_*$. Subtract the exact and approximate
integral equations and apply Gronwall. This gives
$\sup_t\|z-\widehat z\|_*\le E$.

The actual endpoint integrand is
$G(s)=\bar\Psi_{p,x(1,s)}(P_1,P_2,P_3)$. Differentiating in
a unit simplex direction produces one form-coefficient term
bounded by $W_1B_1^4$ and three derivative-vector terms with
total bound $3W_0B_1^2B_2$. Thus $G$ has Lipschitz constant
$C_q$, and the quadrature error on a simplex is at most its
volume times $C_qh$.

At a sample point, first change the point of evaluation of the
form, then change its three arguments one at a time. The
first-jet bound and $\|P\|,\|\widehat P\|\le B_1+1$ give
error at most $C_jE$. Scalar evaluation adds $\xi$.
Summing with exact cell volumes and absolute chain coefficients
gives \eqref{eq:absolute-period-error}.
\end{proof}

An approximate $\widehat P$ need not be the derivative of a
global approximate chain. It approximates the actual first jet
at one integration point, and the displayed estimate compares it
with the actual transported chain.

\subsection{Finite refinement and the observation budget}

We justify that every positive integral tolerance is attained.
First-jet solutions are $C^1$ on the compact time interval and
have $\|P\|\le B_1$. Rational piecewise polynomial curves are
dense in $C^1$ there. Consequently approximations exist with
arbitrarily small initial errors and residuals, and with a strict
margin inside $\|\widehat P\|<B_1+1$.

Enumerate such curves and finite time meshes, dovetailing the
tests. Initial enclosures, polynomial derivatives and the known
variation bounds for $\mathcal F_1$ give upper bounds for the
initial error, residual and derivative-block norm on every mesh
cell. Accept only after the prescribed strict inequalities are
verified. A sufficiently accurate polynomial approximation has
a strict margin, so finite refinement eventually verifies one.
This proves termination of a residual-controlled calculation,
without presuming an error certificate for an unspecified solver.
There are only finitely many sample points at each simplex
subdivision.

For a requested endpoint error $\eta>0$ on a nonzero chain choose
\[
 h<\frac{\eta}{3S\max(1,C_q)},\qquad
 \xi<\frac{\eta}{3S},\qquad
 E<\min\left(\frac12,\frac{\eta}{3S\max(1,C_j)}\right).
\]
The preceding search achieves the required $E$ by taking
$e_0<E/(2e^{\Lambda_*})$ and
$r<E\Lambda_*/[2(e^{\Lambda_*}-1)]$.
All choices can use rational tolerances and upper enclosures of
the constants. Identity transport gives the initial integral.
Evaluate initial and endpoint periods to error $\eta/2$ each;
their difference has absolute error at most $\eta$, irrespective
of cancellation.

The local period is $t_i=p_i$ and is evaluated from the given
coordinate name, or equivalently from the branch-free local
series of Section~\ref{sec:local-periods}. Refining the $1072$
complex jumps and eight local periods supplies all observations
needed by Lemma~\ref{lem:integer-separation}. This proves
Theorem~\ref{thm:accessible-observations}.

Every construction is on the actual compact fibre. The horizontal
field uses an auxiliary ambient metric and equation Jacobians;
it does not require a Ricci-flat metric, a Hermitian--Einstein
metric, or an eigenform evaluator. For every fixed smooth target
and chosen loop its constants are finite and its refinement
terminates. No uniform running time or uniform internal precision
is asserted as the loop approaches the discriminant.
